\documentclass[12pt,a4paper]{article}
% Compile with XeLaTeX or Tectonic. All text is UTF-8.
\usepackage{fontspec}
\setmainfont[Path=fonts/,BoldFont=cmunbx.otf,ItalicFont=cmunti.otf,BoldItalicFont=cmunbi.otf,Ligatures=TeX]{cmunrm.otf}
\usepackage[russian]{babel}
\usepackage[left=25mm,right=20mm,top=20mm,bottom=22mm]{geometry}
\usepackage{amsmath,amssymb,booktabs,array,longtable,graphicx}
\usepackage{float}
\usepackage{hyperref}
\hypersetup{unicode=true,hidelinks,pdftitle={Анализ экспортного потенциала российского натурального меда}}
\urlstyle{same}
\setlength{\parindent}{1.25cm}
\setlength{\parskip}{3pt}
\setlength{\emergencystretch}{3em}
\renewcommand{\arraystretch}{1.18}
\newcommand{\authors}{Сапунова Юлия\\Мамаев Игорь\\Маковский Богдан\\Чеглова Ирина\\Белявцева Виктория}
\begin{document}
\raggedbottom
\hypersetup{pageanchor=false}
\begin{titlepage}
\centering
{\large\textsc{Центральный Университет}\par}
\vspace{1.8cm}
{\Large\bfseries ПРОМЕЖУТОЧНЫЕ РЕЗУЛЬТАТЫ\par ИССЛЕДОВАНИЯ\par}
\vspace{0.5cm}
{\large\bfseries по курсу <<Поиск экспортных рынков>>\par}
\vspace{1.6cm}
\hrule height 1pt
\vspace{0.4cm}
{\LARGE\bfseries Анализ экспортного потенциала и факторов международной торговли\par}
\vspace{0.3cm}
{\Large эмпирическое исследование\par на примере российского натурального меда\par HS2017 040900\par}
\vspace{0.4cm}
\hrule height 1pt
\vspace{1.4cm}
\hfill\begin{minipage}{0.45\textwidth}
\textbf{Выполнили:}\\
\authors\\[0.7cm]
\textbf{Научный руководитель:}\\
Михаил Фролов
\end{minipage}
\vfill
Москва, 2026
\end{titlepage}
\hypersetup{pageanchor=true}
\begingroup
\small
\setlength{\parskip}{0pt}
\tableofcontents
\endgroup
\clearpage
\section*{Аннотация}
В работе сравниваются зарубежные рынки натурального меда (HS 040900), на которых можно подробнее изучить перспективы российских поставок. Использованы данные OEC/BACI за 2018-2024 годы, показатели Всемирного банка и отраслевой обзор Южной Кореи. Российский экспорт вырос с 4,67 до 7,31 млн долларов. В 2024 году на трех крупнейших покупателей пришлось 61,86\% поставок. Сравнение рынков учитывает стоимость и динамику импорта, чистый экспорт и российские поставки. В гравитационной модели торговые потоки положительно связаны с ВВП стран и отрицательно с расстоянием. На тестовой выборке точность модели зависит от выбранного показателя ошибки. Для дальнейшей проверки выбраны США, Ирландия и Южная Корея. США выделяются размером импорта, Ирландия и Южная Корея сочетанием роста импорта с превышением импорта над экспортом. Условия ввоза и расходы на поставку предстоит проверить.

Ключевые слова: натуральный мед, российский экспорт, выбор рынков, RCA, кластеризация, гравитационная модель.

\section{Введение}
Экспортные возможности российского натурального меда нельзя определить только по наличию природных ресурсов или объему зарубежного импорта. В 2024 году российские поставки в сохраненной базе OEC/BACI составили 7,315 млн долларов, а мировая торговля товаром - 2 282,645 млн. Таким образом, доля России равна примерно 0,320\%. При этом более половины российских поставок приходится на три направления. Возникает вопрос, какие рынки заслуживают дополнительного изучения для расширения географии экспорта.

Цель исследования - сопоставить положение России, характеристики импортного спроса и результаты модели торговли и выделить три направления для следующего этапа анализа. Объектом выступает международная торговля натуральным медом. Предметом являются различия импортных рынков и статистические факторы наблюдаемых двусторонних поставок. Код товара HS6 040900 соответствует натуральному меду. Используемый в OEC внутренний идентификатор - 1040900.

Исследование охватывает натуральный мед по коду HS 040900. Торговые данные по этому коду объединяют поставки разных видов меда и форматов упаковки.

\section{Теоретические основы и обзор литературы}
\subsection{Специализация и производство}
Индекс выявленного сравнительного преимущества Балассы сопоставляет долю товара в экспорте страны с долей того же товара в мировом экспорте: RCA = (Xcp / Xc) / (Xwp / Xw). Значение выше единицы указывает на выявленную экспортную специализацию. Показатель характеризует сложившуюся структуру торговли и не измеряет прибыль конкретного предприятия. Поэтому отсутствие специализации не исключает возможностей на отдельных направлениях [2].

Индекс экономической сложности страны ECI и индекс сложности продукта PCI относятся к разным объектам. Они помогают описать структуру специализации, но не измеряют качество меда или прибыльность его продажи. В данном исследовании эти показатели используются как контекст, а не как критерий выбора покупателя.

Ресурсная гипотеза связывает торговлю с ресурсной обеспеченностью. Для пчеловодства важны территории с медоносной растительностью, однако площадь пашни на человека является лишь приближением к этому ресурсу. FAO отмечает возможности пчеловодства на различных типах территорий [3]. Показатель не учитывает состав растений, климат, здоровье пчелиных семей и производительность. Мы ожидаем связь обеспеченности подходящими территориями с экспортом, но не считаем наличие пашни достаточным объяснением.

\subsection{Неоднородность спроса и факторы торговли}
Гравитационная гипотеза относится к размеру рынков и торговым издержкам. Более крупные экономики могут иметь больше возможностей производства и закупок. Расстояние может отражать часть издержек взаимодействия. В гравитационной модели ожидаются положительные коэффициенты при ВВП стран и отрицательный при расстоянии. Эти зависимости проверяются как статистические связи, поскольку производственная специализация, условия доступа и устойчивые торговые отношения учтены неполностью.

Доход на человека характеризует общий уровень доходов, импорт показывает масштаб внешних закупок, а рост импорта отражает изменение их стоимости. Эти показатели не заменяют друг друга. Наличие рецензируемой отраслевой эконометрической статьи пока остается пробелом обзора: имеющиеся материалы дают методологию RCA и отраслевой пример Кореи, но не обосновывают перенос конкретной оценки эластичности спроса на российский мед.

\subsection{Условия доступа и связь с литературой}
В обзоре Балассы [2] определяется смысл выявленной специализации. Публикация FAO [3] уточняет ресурсную гипотезу: пчеловодство использует различные территории, поэтому пашня не охватывает все релевантные природные условия. Методология BACI [4] объясняет согласование зеркальных торговых записей. Обзор AAFC [1] дополняет агрегированную стоимость импорта сведениями о массе и поставщиках. Эти источники отвечают на разные вопросы и не являются взаимозаменяемыми подтверждениями перспективности рынка.

Отраслевой обзор AAFC по данным Global Trade Tracker показывает рост массы корейского импорта в 2024 году на 50,1\% при снижении стоимости на 5,6\%. Средняя стоимость килограмма уменьшилась примерно на 37,1\%. Это может отражать изменение состава поставок. Показатель не является ценой неизменного товара. Новая Зеландия лидирует по стоимости, Вьетнам - по массе. Источник помогает поставить вопрос о подходящем сегменте российского продукта [1].

В обзоре импорт за 2024 год равен 17,176 млн долларов против 16,457 млн в BACI. Разница около 4,37\% относительно BACI. Источники не смешиваются в межстрановой таблице: BACI согласует зеркальную статистику и корректирует различия оценок потоков [4]. Это возможное методологическое объяснение, но точная декомпозиция расхождения не выполнена. Вклад обзора состоит в уточнении сегментации, а не в подтверждении готовой экспортной возможности.

Документы о действующих правилах ввоза для выбранных стран еще не собраны в сопоставимом виде. Возможные расходы на испытания, документы и согласование поставки здесь не оценены. Отсутствие такой оценки ограничивает практический вывод, но не отменяет сравнение спроса на исследовательском этапе.

\subsection{Гипотезы и вклад исследования}
Гипотеза H1: обеспеченность подходящими природными ресурсами связана с относительным экспортом меда. Проверка проводится сопоставлением стран и NEX. Исключения либо слабый ресурсный показатель ограничивают вывод. Гипотеза H2: большие экономики торгуют больше, а удаленность связана с меньшими потоками. Проверка состоит в оценке знаков коэффициентов и качества прогноза. Гипотеза H3: добавление дохода или кластеров улучшает прогноз на общей выборке. Ее поддержка требует устойчивого снижения ошибок, а не только ожидаемого знака коэффициента.

Для выбора рынков мы сопоставили размер импорта, его изменение за три года и поставки из России.

\section{Данные и методы исследования}
\subsection{Источники, единицы наблюдения и периоды}
Основой служит сохраненная 4 ноября 2026 года выгрузка OEC из базы trade\_i\_baci\_a\_17, HS2017, 2018-2024 годы [5]. Она содержит 25 219 зарегистрированных двусторонних потоков. Стоимость выражена в текущих долларах США. Для таблиц и рисунков используются миллионы долларов. Мировой экспорт определяется суммой двусторонних потоков без повторного сложения экспорта и импорта.

Описательные итоги рассчитаны по исходной торговой выгрузке. Для регрессии нужны также ВВП, доход на человека и координаты столиц из World Bank [6]. После проверки наличия признаков осталось 24 152 наблюдения. Исключение 1 067 строк затрагивает около 1,52\% стоимости торговли. Поэтому агрегаты очищенной панели не подменяют полные торговые итоги. Отсутствующая строка обозначается как отсутствие зарегистрированного наблюдения, а не автоматически как нулевой поток.

Для сравнения спроса выбран общий период 2021-2024 годов. Изменение стоимости равно (M2024 / M2021 - 1) × 100\%. Среднегодовой темп - ((M2024 / M2021)\textsuperscript{1/3} - 1) × 100\%. Нормированный чистый экспорт NEX = (X - M) / (X + M). Отрицательное значение означает превышение импорта над экспортом. Оно не измеряет дефицит внутреннего производства относительно потребления. CR3 равен сумме долей трех крупнейших поставщиков.

Показатели RCA, ECI, PCI и состав экспортной корзины взяты из материалов недели 4. Для анализа потоков используется выгрузка OEC/BACI из работы недели 8.

\subsection{Показатели предложения и группировка рынков}
\begin{align}
RCA_{cp}&=\frac{X_{cp}/X_c}{X_{wp}/X_w}, &
NEX_c&=\frac{X_c-M_c}{X_c+M_c},\\
HHI&=\sum_{k=1}^{K}s_k^2, &
H&=-\sum_{k:s_k>0}s_k\ln s_k,\\
RX_c&=\frac{X_{c,040900}}{X_{c,190590}}.
\end{align}
В формуле RCA индексы c и p обозначают страну и товар, w обозначает мир. В формуле NEX экспорт и импорт относятся к меду. Для HHI и энтропии s обозначает долю группы или покупателя. Объект концентрации каждый раз уточняется. В RX товар сравнения представляет собой прочие хлебобулочные изделия HS2017 190590. Значения RCA ниже единицы не означают невозможность прибыльной поставки.

Кластеризация и прогнозная проверка воспроизводят выполненную работу недели 8. Модели обучаются на 2018-2022 годах и проверяются на 2023-2024 годах. Тридцать рынков отбираются по данным обучения. Для кластеризации стандартизируются логарифм импорта, его среднегодовой рост и ВВП на человека. Три группы построены без использования тестовых лет. В описательном выборе рынков данные 2024 года используются открыто: это текущий аналитический отбор, а не ретроспективная проверка инвестиционного решения.

Использован алгоритм K-means с евклидовым расстоянием после стандартизации. Число групп равно трем по условию работы. Наличие группы из одной страны учитывается при интерпретации. 

\subsection{Гравитационная и панельные модели}
\begin{equation}
\ln T_{ijt}=\alpha+\beta_1\ln GDP_{it}+\beta_2\ln GDP_{jt}
+\beta_3\ln d_{ij}+u_{ijt}.
\label{eq:gravity}
\end{equation}
T обозначает положительный поток в млн текущих долларов, GDP обозначает ВВП в текущих долларах, d обозначает расстояние между столицами в километрах. Первое расширение добавляет логарифм ВВП на человека импортера, второе добавляет индикаторы кластеров. Первая группа служит базой сравнения. Обратное преобразование выполняется через экспоненту без дополнительной поправки, как в задании.

Регрессия оценивается для логарифмов положительных потоков. Качество проверяется по MAE, RMSE и R\textsuperscript{2} с сопоставлением с постоянными прогнозами. Для расширений базовая модель переоценивается на общей выборке. Полученные результаты описывают имеющуюся торговлю и не моделируют возникновение новых связей.

\begin{align}
MAE&=\frac{1}{n}\sum_{k=1}^{n}|T_k-\widehat T_k|,\\
RMSE&=\sqrt{\frac{1}{n}\sum_{k=1}^{n}(T_k-\widehat T_k)^2}.
\end{align}
R-квадрат рассчитывается отдельно в уровнях и логарифмах. Постоянные прогнозы по среднему и медиане используют только обучение. Расширения сравниваются с базовой моделью на одних и тех же 7 589 обучающих и 3 190 тестовых наблюдениях. Оценки pooled OLS, FE и RE для российской панели относятся к следующей стадии курса и в статью как выполненные результаты не включены.

\subsection{Сравнение условий доступа}
Для сравнения условий доступа нужно уточнить требования к ввозу российского меда и рассчитать расходы для одинаковой партии по каждой стране.

\section{Анализ структуры и динамики мирового экспорта выбранного товара}
\subsection{Ключевые игроки на мировом рынке}
Таблица 1 показывает различия мировой и российской динамики. В 2024 году мировая стоимость торговли была на 15,73\% ниже 2021 года, а российская - на 5,17\% выше. За 2018-2024 годы российский экспорт вырос примерно на 56,7\%, хотя внутри периода изменения были неравномерными. Стоимостной рост нельзя автоматически интерпретировать как рост физических поставок.

\begin{table}[H]
\centering
\caption{Мировая торговля натуральным медом и экспорт России}\label{tab:years}
{\small
\begin{tabular}{lrr}
\toprule
Год & Мировая торговля млн долл. & Россия млн долл. \\ \midrule
2018 & 2 245,40 & 4,668 \\
2019 & 1 991,67 & 5,540 \\
2020 & 2 265,61 & 5,178 \\
2021 & 2 708,92 & 6,955 \\
2022 & 2 739,97 & 5,456 \\
2023 & 2 187,67 & 6,100 \\
2024 & 2 282,65 & 7,315 \\
\bottomrule
\end{tabular}}
\par\smallskip{\footnotesize Источник: расчеты по сохраненной выгрузке OEC/BACI, HS2017 040900.}
\end{table}
\begin{table}[H]
\centering
\caption{Крупнейшие экспортеры натурального меда в 2024 году}\label{tab:exporters}
{\small
\begin{tabular}{lrr}
\toprule
Экспортер & Экспорт млн долл. & Доля мира проц. \\ \midrule
Новая Зеландия & 251,31 & 11,01 \\
Китай & 239,01 & 10,47 \\
Аргентина & 190,37 & 8,34 \\
Индия & 187,13 & 8,20 \\
Украина & 167,46 & 7,34 \\
Германия & 132,80 & 5,82 \\
Испания & 101,96 & 4,47 \\
Бразилия & 101,11 & 4,43 \\
\bottomrule
\end{tabular}}
\par\smallskip{\footnotesize Источник: расчеты по сохраненной выгрузке OEC/BACI, HS2017 040900.}
\end{table}
Таблица с ведущими экспортерами показывает международную структуру предложения. Доля России значительно меньше долей лидеров. Различия торговой стоимости не позволяют непосредственно сравнить качество, физический выпуск или прибыль производителей.

\subsection{Выявленное сравнительное преимущество (RCA)}
По таблице OEC из работы недели 4 RCA российского меда за 2024 год составил 0,195. Доля меда в российском экспорте ниже его доли в мировом экспорте. Значение приведено по той же выгрузке, которую использовали в работе недели 4.

\subsection{Концентрация и сложность экспорта}
По выгрузке из работы недели 4 индекс концентрации HHI товарной корзины России на уровне HS2 составляет около 0,399. На крупнейшую товарную группу приходится 62,5\% экспорта в этой выгрузке.

Показатели ECI России 0,330 и PCI меда -1,027 перенесены из архивных результатов работы недели 4. Сочетание значений не позволяет сделать вывод о качестве или прибыльности продукта. HHI всей товарной корзины отличается от концентрации направлений экспорта меда, рассчитанной в разделе 5.

\subsection{Обеспеченность ресурсами и структура торговли}
В выполненной ресурсной проверке показатель пашни на человека в России равен 0,8467 га, в Японии - 0,03245 га, то есть различается примерно в 26,1 раза. Отношение экспорта меда к экспорту товара сравнения HS2017 190590 составляет соответственно 0,0340 и 0,00174. Порядок согласуется с ресурсной гипотезой, если предположение об относительной землеемкости меда верно. Но ресурсная интенсивность производства непосредственно не измерена.

В выборке из 16 стран доля чистых экспортеров составляет 50\% среди стран выше медианы обеспеченности пашней и 25\% ниже нее. Одновременно наблюдаются исключения: США импортируют больше, чем экспортируют, при сравнительно высокой обеспеченности пашней. Новая Зеландия является чистым экспортером при низком значении выбранного показателя. Это ограничивает объясняющую роль пашни и указывает на необходимость учитывать производительность, природную среду и специализацию. Эти результаты перенесены из выполненной работы недели 5, а не получены повторной загрузкой WDI.

\subsection{Сравнение импортных рынков и кластеров}
Таблица \ref{tab:markets} сопоставляет шесть ранее изученных импортеров. Все значения торговли повторно рассчитаны на полной сохраненной выгрузке. Российские поставки приведены в тысячах долларов. Н/д означает отсутствие строки в исходной базе за этот год.

\begin{table}[H]
\centering
\caption{Импортные рынки натурального меда}\label{tab:markets}
{\small
\begin{tabular}{lrrrr}
\toprule
Рынок & Импорт млн долл. & Рост проц. & NEX & Из РФ тыс. долл. \\ \midrule
США & 621,24 & -0,84 & -0,937 & 201,869 \\
Германия & 248,97 & -26,33 & -0,304 & 1,679 \\
Нидерланды & 74,51 & 10,54 & -0,562 & н/д \\
Ирландия & 24,81 & 40,26 & -0,710 & 0,093 \\
Япония & 132,24 & -21,58 & -0,991 & н/д \\
Южная Корея & 16,46 & 9,44 & -0,976 & 5,522 \\
\bottomrule
\end{tabular}}
\par\smallskip{\footnotesize Источник: расчеты по сохраненной выгрузке OEC/BACI, HS2017 040900.}
\end{table}
Импорт и российские поставки относятся к 2024 году, рост рассчитан за весь период 2021-2024 годов. Обозначение н/д означает отсутствие строки российской поставки. Сравнение приведено в таблице \ref{tab:markets}.

Самый крупный рынок - США, однако его стоимость импорта почти не изменилась относительно 2021 года. Германия и Япония велики по масштабу, но демонстрируют сокращение стоимости. Ирландия растет быстрее остальных пяти, имея значительно меньший импорт. Корея сочетает умеренный стоимостной рост и почти предельное отрицательное NEX. Нидерланды крупнее Ирландии и Кореи, но российская поставка за 2024 год в исходной выгрузке отсутствует.

В кластеризации недели 8 первая группа включает 20 рынков со средним импортом 110,62 млн долларов и доходом 57,84 тыс. долларов на человека. Третья включает девять рынков со средним импортом 28,42 млн и более высоким средним ростом 11,31\% в год. США и Ирландия относятся к первой группе, Корея - к третьей. Эти средние относятся к периоду и панели кластеризации, а не к описательным итогам 2024 года. Вторая группа состоит только из Нигерии с очень высоким темпом роста с низкой базы. Ее нельзя считать устойчивым самостоятельным типом рынков. Число групп задано условием учебной работы.

\begin{figure}[H]\centering
\includegraphics[width=\textwidth]{markets.png}
\caption{Размер и изменение стоимости импорта шести рынков}
\label{fig:markets}
\par\smallskip{\footnotesize Источник: расчеты по OEC/BACI. Импорт 2024 года, рост за период 2021-2024 годов.}
\end{figure}
\section{Россия на мировом рынке выбранного товара}
\subsection{Структура и динамика российского экспорта}
По таблице динамики российский экспорт вырос с 4,668 млн долларов в 2018 году до 7,315 млн в 2024 году. Ряд не является монотонным: сокращение наблюдалось в 2020 и 2022 годах. Для описания используется одна база и один период, поэтому ранний ряд 2015-2025 годов с неподтвержденными различиями источников не переносится. Доля России в мировой торговле 2024 года равна 0,320\%.

\begin{figure}[H]\centering
\includegraphics[width=\textwidth]{world.png}
\caption{Мировая торговля и российский экспорт натурального меда}
\label{fig:world}
\par\smallskip{\footnotesize Источник: расчеты по OEC/BACI, 2018-2024 годы, млн текущих долларов.}
\end{figure}
На рисунке показана различная динамика мирового рынка и России. Концентрация российских покупателей, рассчитанная на положительных потоках 2024 года, составляет HHI 0,1563. Это отдельный показатель, который не смешивается с HHI товарной корзины по HS2.

\subsection{География поставок и выбор направлений для сравнения}
Крупнейшими покупателями российского меда в 2024 году были Словакия - 1,947 млн долларов, Сербия - 1,346 млн и Китай - 1,232 млн. Их совокупная доля равна 61,86\%, доля первого направления - 26,62\%, первых пяти - 78,58\%. Ранее приведенные в первой домашней работе 36,3\%, 70\% и 82\% не переносятся: они не воспроизводятся на единой выгрузке 2024 года. Концентрация поддерживает постановку вопроса о диверсификации, но сама по себе не доказывает, что новые рынки будут выгоднее существующих.

\begin{table}[H]
\centering
\caption{Основные покупатели российского меда в 2024 году}\label{tab:dest}
{\small
\begin{tabular}{lrr}
\toprule
Покупатель & Поставки млн долл. & Доля проц. \\ \midrule
Словакия & 1,947 & 26,62 \\
Сербия & 1,346 & 18,41 \\
Китай & 1,232 & 16,84 \\
Казахстан & 0,668 & 9,14 \\
Узбекистан & 0,554 & 7,57 \\
Литва & 0,401 & 5,49 \\
Польша & 0,382 & 5,22 \\
\bottomrule
\end{tabular}}
\par\smallskip{\footnotesize Источник: расчеты по сохраненной выгрузке OEC/BACI, HS2017 040900.}
\end{table}
Из шести рынков, рассмотренных в работе недели 5, для дальнейшей проверки выбраны США, Ирландия и Южная Корея.

США. Импорт равен 621,24 млн долларов, NEX равен -0,937. Российские поставки присутствуют во всех семи годах и составляют 201,9 тыс. долларов в 2024 году. Это основание изучать возможность развития существующей связи. Возражения - почти неизменная стоимость импорта с 2021 года, снижение российских поставок относительно 2023 года и сильные действующие конкуренты. Индия, Аргентина и Бразилия обеспечивают 56,58\% импорта. Прежде чем сохранять страну в финальной рекомендации, необходимо проверить применяемые условия торговли и экономику доступного канала поставки.

Ирландия. Импорт вырос с 17,69 до 24,81 млн долларов, или на 40,26\% за три года, при NEX -0,710. Это сильнейшая динамика среди шести ранее рассматривавшихся рынков. Главный недостаток - крайне малая российская поставка: всего 93 доллара в 2024 году. Наличие такой записи не подтверждает коммерчески действующий канал. Три крупнейших поставщика обеспечивают 74,14\% импорта. Следует проверить происхождение роста, физические объемы и возможную роль перепродаж. При отсутствии подтверждения жизнеспособного канала рынок должен уступить место альтернативе.

Южная Корея. Импорт равен 16,46 млн долларов, изменение к 2021 году составляет +9,44\%, NEX равен -0,976. Российские потоки наблюдаются во всех семи годах, но снизились с 26,4 тыс. долларов в 2023 году до 5,5 тыс. в 2024 году. По BACI на Новую Зеландию, США и Австралию приходится 77,67\% стоимости импорта. Требуется определить целевой продуктовый сегмент и проверить причины сокращения российских поставок.

\section{Эмпирическая оценка факторов и интерпретация результатов}
\subsection{Базовая гравитационная модель}
\begin{align}
\widehat{\ln T}_{ijt}={}&-24{,}9285+0{,}5196\ln GDP_{it}\\
&+0{,}4988\ln GDP_{jt}-0{,}8393\ln d_{ij}.\notag
\end{align}
Знаки коэффициентов соответствуют гипотезе H2. При прочих включенных признаках увеличение ВВП экспортера на 1\% связано с ростом потока на 0,52\%, ВВП импортера на 0,50\%, расстояния со снижением на 0,84\%. Эти оценки не являются причинными эффектами.

\begin{table}[H]
\centering
\caption{Качество прогноза на тесте 2023-2024 годов}\label{tab:model}
{\small
\begin{tabular}{lrrr}
\toprule
Прогноз & MAE млн долл. & RMSE млн долл. & R2 уровни \\ \midrule
Среднее & 1,0765 & 4,4285 & -0,0003 \\
Медиана & 0,6190 & 4,4700 & -0,0191 \\
Гравитационная & 0,6138 & 4,4591 & -0,0141 \\
\bottomrule
\end{tabular}}
\par\smallskip{\footnotesize Источник: сохраненные расчеты работы недели 8.}
\end{table}
На 7 091 тестовом наблюдении MAE регрессии составляет 0,6138 млн долларов против 0,6190 у медианного прогноза. По RMSE модель уступает прогнозу по среднему: 4,4591 против 4,4285 млн. Тестовый R\textsuperscript{2} в логарифмах равен 0,1344, в уровнях - -0,0141. Следовательно, размеры экономик и расстояние передают часть различий потоков, но недостаточны для точного прогноза стоимости. Оценивание логарифмов и обратное преобразование через экспоненту без поправки также влияют на результат в денежных уровнях.

\subsection{Доход и индикаторы групп}
На общей выборке тридцати импортеров добавление дохода уменьшает MAE с 1,243541 до 1,243365 млн долларов. Добавление кластеров - до 1,241052 млн. Улучшения малы, статистическая значимость различий прогнозных ошибок не проверялась. При этом RMSE модели с кластерами повышается в двух из трех групп. Поэтому сегментация полезна для описания рынков, но не обеспечивает существенного общего улучшения прогноза.

\begin{table}[H]
\centering
\caption{Сравнение расширений на общей выборке}\label{tab:model_ext}
{\small
\begin{tabular}{lrrr}
\toprule
Модель & MAE & RMSE & R2 логарифмы \\ \midrule
Базовая & 1,243541 & 6,620579 & 0,147979 \\
С доходом & 1,243365 & 6,620179 & 0,148021 \\
С кластерами & 1,241052 & 6,615725 & 0,153620 \\
\bottomrule
\end{tabular}}
\par\smallskip{\footnotesize Источник: сохраненные расчеты работы недели 8.}
\end{table}
MAE и RMSE представлены в млн долларов. Гипотеза H3 получает только слабую поддержку: некоторые метрики улучшаются, но изменения малы и неодинаковы по группам. Наличие статистически значимого или практически существенного выигрыша в прогнозе не установлено.

\subsection{Российская панель: pooled OLS, FE и RE}
Раздел относится к следующему этапу исследования. По итогам первых восьми недель оценки pooled OLS, FE и RE не получены. В текущей работе представлен российский срез мировой модели, который не заменяет панельную оценку. На следующем этапе потребуется определить состав панели, эффекты стран и лет, способ оценки стандартных ошибок и сопоставимость моделей.

\subsection{Модельные разрывы и основания дальнейшего изучения}
Для российских поставок в Корею за 2024 год модельный прогноз равен 0,0317 млн долларов, факт - 0,0055 млн. Этот разрыв может указывать на неучтенные ограничения или особенности торговли. Он не является оценкой возможной дополнительной выручки. Кроме того, тестовая модель использует фактический ВВП тестовых лет и представляет условный прогноз при известных макропоказателях.

\subsection{Продукт и условия поставки}
Товар исследования определяется кодом HS 040900. Параметры партии и требования к поставке в США, Ирландию и Южную Корею предстоит уточнить.

\subsection{Расходы и объем поставки}
Для расчета расходов нужны цена закупки, характеристики партии, стоимость перевозки и обязательные платежи. Эти данные будут собраны на следующем этапе.

\subsection{Сроки, доля в импорте и проверка гипотезы}
Российская доля в импорте США составляет около 0,0325\%, Ирландии 0,00037\%, Южной Кореи 0,0336\% в 2024 году. Малые доли не измеряют незанятый рынок. Проверка следующего этапа должна сопоставить срок выхода и стоимость конкретной партии с существующими поставщиками. Гипотеза привлекательности направления пересматривается, если доступ отсутствует, целевой сегмент не растет или сопоставимая поставка не выдерживает конкуренции по цене и требованиям.

\subsection{Обобщающие выводы и рекомендации}
При сравнении трех стран нужно проверить, можно ли нарастить существующие поставки в США и Южную Корею и есть ли в Ирландии покупатели российского меда.

\section{Заключение}
Российский натуральный мед занимает небольшую долю мировой торговли и зависит от ограниченного числа покупателей. Ресурсная обеспеченность и общие макроэкономические признаки помогают формулировать гипотезы, но недостаточны для выбора рынка. Гравитационная модель дает ожидаемые знаки коэффициентов, одновременно демонстрируя ограниченную точность в денежном выражении.

Из рассмотренных стран США выделяются размером рынка, Ирландия ростом стоимости импорта, а Южная Корея сочетанием роста импорта и истории российских поставок. Сравнение условий ввоза и расходов на поставку позволит уточнить этот выбор.

\section{Список литературы, оформление и приложения}
\subsection{Научные и аналитические источники}
\noindent [1] Agriculture and Agri-Food Canada. Customized Report Service - Sweet spreads and honey in South Korea. 2025. \url{https://agriculture.canada.ca/en/international-trade/reports-and-guides/customized-report-service-sweet-spreads-and-honey-south-korea}. Проверено 09.11.2026. Таблицы Global Trade Tracker, 2025.

\noindent [2] Balassa B. Trade Liberalisation and Revealed Comparative Advantage. The Manchester School. 1965. Vol. 33, № 2. P. 99-123. \url{https://doi.org/10.1111/j.1467-9957.1965.tb00050.x}. Использован разбор статьи из материалов команды.

\noindent [3] Bradbear N. Beekeeping and sustainable livelihoods. FAO Diversification booklet 1. 2004. \url{https://www.fao.org/4/y5110e/y5110e02.htm}. Страница проверена 09.11.2026.

\noindent [4] Gaulier G., Zignago S. BACI International Trade Database at the Product-Level The 1994-2007 Version. CEPII Working Paper 2010-23. \url{https://www.cepii.fr/PDF_PUB/wp/2010/wp2010-23.pdf}. 

\noindent [5] OEC / CEPII BACI. HS6 REV. 2017, 2018-2024. \url{https://api.oec.world/tesseract/cubes/trade_i_baci_a_17}. Торговая выгрузка от 04.11.2026, 25 219 строк. Данные: \url{https://api.oec.world/tesseract/data.jsonrecords}.

\noindent [6] World Bank. World Development Indicators. GDP current US\$, NY.GDP.MKTP.CD. GDP per capita current US\$, NY.GDP.PCAP.CD. Arable land hectares per person, AG.LND.ARBL.HA.PC. \url{https://data.worldbank.org/}. Макропоказатели модели выгружены 04.11.2026. Земельный показатель перенесен из выполненной работы недели 5.

\subsection{Таблицы, рисунки и язык изложения}
В таблицах и на рисунках торговые потоки приведены в текущих долларах. Обозначение «н/д» означает, что в выгрузке отсутствует запись о поставке.

\subsection{Рефлексия и вклад команды}
При объединении работ мы сверили исходные данные и заново рассчитали доли крупнейших покупателей российского меда по выгрузке за 2024 год. Индивидуальный вклад участников нужно уточнить перед сдачей.

\end{document}
